\section{Building The Corpus}

\subsection{Alcide De Gasperi's Corpus}

I would like to thank Sara Tonelli for providing me with the De Gasperi's Corpus \cite{tonelli_prendo_2019}, a collection of Alcide De Gasperi's public documents with gold and silver annotation.

The corpus is a collection of 2,762 documents issued between 1901 and 1954, formatted into XML files which include metadata that covers not only the title, the date and the place of publication, but also key-concepts automatically extracted from each text (with the corresponding relevance score) and genre labels manually assigned by domain experts. Furthermore, the release includes silver annotation for lemma, part of speech, person names and place names with associated coordinates in a CoNLL-like format.

An example of the XML formatting can be seen at listing \ref{code:xml_degasperi}. 
Thanks to the \textit{<genres>} tag, I was able to extract 474 public speeches, each of which with a location, a precise date and a list of keywords.

\begin{lstlisting}[language={XML},label={code:xml_degasperi}, caption={An example of how a single speech in De Gasperi's Corpus is formatted in XML}]
<?xml version="1.0" encoding="UTF-8" standalone="no"?>
<document id="I.doc_7">
   <publication_date>1901-11-26</publication_date>
   <publication_place>
      <place_name>Trento</place_name>
      <place_latitude>46.0664228</place_latitude>
      <place_longitude>11.1257601</place_longitude>
   </publication_place>
   <keywords>
      <keyword score="39.16">first keyword</keyword>
      ...
      <keyword score="5.59">last keyword</keyword>
   </keywords>
   <genres>
      <genre>Speech Type</genre>
   </genres>
   <text>The speech itself if stored between these brackets</text>
</document>
\end{lstlisting}

\subsection{Giorgia Meloni's Corpus}

My objective was to build Meloni's corpus by transcribing videos of her public speeches, such as talks and interviews. Fortunately, there is an unofficial YouTube channel (that I reached from the Prime Minister's official website) which aggregates more than four thousand videos of her public appearances. The channel is called "Giorgia Meloni News"\footnote{Giorgia Meloni News: \url{https://www.youtube.com/@GiorgiaMeloniTv}}.

Given a YouTube URL, I can use Python's library \textit{yt-dlp} to retrieve the video's metadata and download its audio content in mp3 format, as shown in listing \ref{code:yt_download}. I had to manually pass the cookies taken from my web-browser, and I used some extra commands to prevent YouTube to block the requests due to suspicious activity. 

Then, the mp3 file just retrieved gets injected into OpenAI's Whisper library, which uses an encoder-decoder Transformer to transcribe it into text, as shown in listing \ref{code:whisper}. The model I used is the \textit{"medium"}\footnote{Whisper model sizes available are tiny, base, medium and large.}, which at 5 GB fits within the total 6 GB of VRAM available to my RTX 4050 GPU. The model was instructed to predict Italian. 

The difference in precision between the models "tiny" (which requires less than 1 GB of VRAM) and "medium" is quite high, as can be seen in the following pieces of text that have been generated from the same audio file:

\begin{itemize}
    \item \textbf{Model tiny:} "iamo con se è beruto fuori da questa rio neone sotto il profiro tecnico per quanto riguarda la franha e come il governo un tino era ad essere vicino alla popolazione di Nishenia."
    \item \textbf{Model medium:} "Diciamo cosa è venuto fuori da questa riunione sotto il profilo tecnico per quanto riguarda la frana e come il governo continuerà ad essere vicino alla popolazione diniscemeca."
\end{itemize}

The last step of the pipeline saves the transcription into a csv file, along with the extracted metadata. Each file has the following fields: \textit{politician} ("meloni" in this case), \textit{historical\_date} (the upload date of the video), \textit{location} and \textit{tags} (extracted from the metadata if available, empty strings otherwise), \textit{description} and \textit{title} of the video, \textit{url} which stores the permalink of the video itself, and lastly \textit{text} which holds the whole transcription.

Each video is processed immediately after being retrieved, and until it has been saved in a csv file the program does not try to fetch the next one. This is done for two reasons: first, to prevent \textit{yt-dlp} from sending requests too close to each other, which could result in YouTube denying them due to suspicious activity. Second, even if the pipeline halts for any reason, the videos processed up to that point will not be lost.  

\begin{lstlisting}[language={python},label={code:yt_download}, caption={Downloading audio and metadata from a YouTube video}]
result = subprocess.run([
    "yt-dlp",
    "--print-json",
    "-x",                               # download audio only
    "--audio-format", "mp3",            
    "--cookies", "yt_cookies.txt",      # manual cookies, extract with browser extension
    "--sleep-requests", "2",            # Sleep 2s between requests
    "--sleep-interval", "5",            # Sleep 5s between downloads
    "--max-sleep-interval", "15",       # Randomize up to 15s
    "--limit-rate", "5M",               # Throttle to 5MB/s (mimics streaming)
    "-o", audio_file.replace('.mp3', ''), 
    video_url
], capture_output=True, text=True, check=True)
\end{lstlisting}

\begin{lstlisting}[language={python},label={code:whisper}, caption={Transcribe the audio file into text using Whisper}]
subprocess.run([
    "whisper",
    "--language",        lang_code,     # "Italian" chosen
    "--word_timestamps", "True",
    "--model",           model_name,    # model "medium" chosen
    "--output_dir",      output_dir,
    "--device",          "cuda",        # ensures the model gets loaded in GPU
    audio_file
], check=True)
\end{lstlisting}

\subsection{Corpora Description}

Table \ref{tab:datasets_summary} describes the corpora collected so far: the date each corpus was compiled, the number of speeches contained, the period in which the speeches were originally delivered (represented as a date range, from the oldest to the most recent speech), and the median speech length in words

\begin{table}[h]
    \centering
    \caption{Summary of Meloni and De Gasperi corpora.}
    \label{tab:datasets_summary}
    \begin{tabular}{lcc}
        \toprule
        \textbf{Metric} & \textbf{Meloni Corpus} & \textbf{De Gasperi Corpus} \\
        \midrule
        Speeches Date Range           & 2025-05-16 $\rightarrow$ 2026-05-21 & 1901-09-21 $\rightarrow$ 1954-07-24 \\
        Number of Speeches             & 72                                  & 473                                 \\
        Median Words/Speech  & 982                                 & 701                                 \\
        Compiled On          & May 2026                            & July 2019                           \\
        \bottomrule
    \end{tabular}
\end{table}

